\documentclass[12pt]{article}

% Basic packages
\usepackage{amsmath,amssymb}
\usepackage{graphicx}
\usepackage{booktabs}
\usepackage[margin=1in]{geometry}
\usepackage{float}
\usepackage{xcolor}

% Define colors for code
\definecolor{codegreen}{rgb}{0,0.6,0}
\definecolor{codegray}{rgb}{0.5,0.5,0.5}
\definecolor{codepurple}{rgb}{0.58,0,0.82}
\definecolor{backcolour}{rgb}{0.95,0.95,0.92}

% Title
\title{\textbf{Neural ODEs for Langmuir Adsorption Kinetics}}
\author{Kasinadh Sarma}
\date{May 2, 2025}

% Load hyperref last
\usepackage{hyperref}
\hypersetup{
    colorlinks = true,
    linkcolor = blue,
    urlcolor  = blue,
    citecolor = blue,
}

\begin{document}

\pagenumbering{roman}

\maketitle
\clearpage

\begin{abstract}
This report presents a comprehensive implementation and analysis of Neural Ordinary Differential Equations (Neural ODEs) applied to modeling Langmuir adsorption kinetics. We compare two approaches: traditional parameter estimation and Neural ODEs, evaluating their performance in predicting adsorption dynamics. Our results demonstrate that Neural ODEs can effectively capture the underlying dynamics while offering flexibility in modeling complex systems.
\end{abstract}
\clearpage

\pagenumbering{arabic}

\tableofcontents
\clearpage

\section{Introduction}

\subsection{Background}
Langmuir adsorption is a fundamental model in surface chemistry describing the adsorption of molecules onto solid surfaces. This model, developed by Irving Langmuir in 1918, has become a cornerstone in understanding gas-solid and liquid-solid interface phenomena. Its applications span across various fields including:

\begin{itemize}
    \item Catalysis and heterogeneous reactions
    \item Gas storage and separation
    \item Water purification and treatment
    \item Surface chemistry and materials science
\end{itemize}

The model makes several key assumptions:
\begin{itemize}
    \item Fixed number of adsorption sites that are all equivalent
    \item Monolayer coverage (only one layer of adsorbate molecules)
    \item Energetically equivalent sites (uniform surface)
    \item No interactions between adjacent adsorbed molecules
    \item Dynamic equilibrium between adsorption and desorption processes
\end{itemize}

\subsection{Mathematical Framework}
The Langmuir adsorption model is described by a first-order ordinary differential equation:

\begin{equation}
    \frac{d\theta}{dt} = k_a P (1-\theta) - k_d \theta
\end{equation}

where:
\begin{itemize}
    \item $\theta$: fractional surface coverage ($0 \leq \theta \leq 1$)
    \item $k_a$: adsorption rate constant (kinetic parameter)
    \item $k_d$: desorption rate constant (kinetic parameter)
    \item $P$: pressure/concentration of adsorbate
\end{itemize}

At equilibrium ($\frac{d\theta}{dt} = 0$), this leads to the Langmuir isotherm:

\begin{equation}
    \theta_{eq} = \frac{K_L P}{1 + K_L P}
\end{equation}

where $K_L = k_a/k_d$ is the Langmuir equilibrium constant.

\section{Methodology}

\subsection{Implementation Details}
Our implementation leverages Julia's scientific computing ecosystem:

\subsubsection{Core Libraries}
\begin{itemize}
    \item \texttt{DifferentialEquations.jl}: Robust ODE solvers with adaptive time-stepping
    \item \texttt{Flux.jl}: Deep learning framework for neural network architecture
    \item \texttt{DiffEqFlux.jl}: Integrates differential equations with neural networks
    \item \texttt{Optimization.jl}: Advanced optimization algorithms
\end{itemize}

\subsubsection{Data Generation}
\begin{itemize}
    \item Synthetic data generated using known parameters ($k_a = 0.5$, $k_d = 0.1$)
    \item Time range: 0 to 20 time units with 100 equally spaced points
    \item Added Gaussian noise ($\sigma = 0.02$) to simulate experimental conditions
    \item 80-20 split for training and validation sets
\end{itemize}

\subsection{Neural ODE Architecture}

\subsubsection{Network Structure}
\begin{itemize}
    \item Input layer: 1 unit (current state $\theta$)
    \item Hidden layer: 16 units with tanh activation
    \item Output layer: 1 unit ($d\theta/dt$ prediction)
\end{itemize}

\subsubsection{Training Configuration}
\begin{itemize}
    \item Optimizer: ADAM with learning rate = 0.01
    \item Loss function: Mean squared error
    \item Batch size: 32 samples
    \item Training epochs: 1000 with early stopping
    \item Regularization: L2 with coefficient $\lambda = 0.01$
\end{itemize}

\section{Results and Discussion}

\subsection{Model Performance}

\begin{table}[htbp]
\centering
\begin{tabular}{@{}lcccc@{}}
\toprule
Model & RMSE (True) & RMSE (Noisy) & R² (True) & MAE (True) \\
\midrule
Neural ODE & 0.03976 & 0.04248 & 0.94963 & 0.02794 \\
Parameter Estimation & 0.00217 & 0.00891 & 0.99985 & 0.00192 \\
\bottomrule
\end{tabular}
\caption{Performance comparison between Neural ODE and Parameter Estimation approaches}
\label{tab:performance}
\end{table}

\subsection{Visual Comparisons}

\begin{figure}[htbp]
    \centering
    \includegraphics[width=0.8\textwidth]{langmuir_comparison.png}
    \caption{Comparison between true solution (solid line), Neural ODE predictions (dashed), and parameter estimation (dotted) for clean and noisy data}
    \label{fig:comparison}
\end{figure}

\begin{figure}[htbp]
    \centering
    \includegraphics[width=0.8\textwidth]{langmuir_neural_ode_result.png}
    \caption{Neural ODE model predictions showing uncertainty bounds and training points}
    \label{fig:neural_ode_results}
\end{figure}

\subsection{Discussion}

\subsubsection{Model Accuracy}
\begin{itemize}
    \item Parameter estimation shows superior accuracy (RMSE = 0.00217)
    \item Neural ODEs demonstrate robust performance with noisy data
    \item Both approaches capture essential dynamics
\end{itemize}

\subsubsection{Computational Considerations}
\begin{itemize}
    \item Neural ODEs: longer training ($\sim$2.5s) but more flexible
    \item Parameter estimation: faster ($\sim$0.5s) but needs good initial guesses
    \item Neural ODEs scale better with system complexity
\end{itemize}

\section{Conclusion}

Our comprehensive study demonstrates that while parameter estimation shows superior accuracy for this well-defined system, Neural ODEs offer valuable advantages in terms of flexibility and robustness to noise. The methodology developed here provides a framework for future studies in chemical kinetics modeling.

\section{Future Work}

\begin{itemize}
    \item Testing with more complex adsorption models
    \item Exploring different neural network architectures
    \item Investigating noise impact on model performance
    \item Extending to multi-component systems
    \item Integration with real-time control systems
\end{itemize}

\clearpage
\begin{thebibliography}{9}
\bibitem{chen2018} Chen, R. T., et al. (2018). Neural Ordinary Differential Equations. arXiv:1806.07366.
\bibitem{rackauckas2019} Rackauckas, C., et al. (2019). DiffEqFlux.jl - A Julia Library for Neural Differential Equations.
\bibitem{langmuir1918} Langmuir, I. (1918). The adsorption of gases on plane surfaces of glass, mica and platinum.
\end{thebibliography}

\end{document}